\documentclass[12pt]{article}
\usepackage[T1]{fontenc}
\usepackage[utf8]{inputenc}
\usepackage{lmodern}
\usepackage{textcomp}
\usepackage{enumitem}
\usepackage{geometry}
\geometry{margin=1in}
\begin{document}
\begin{center}
Answer Key - Chemistry - Graduate Level Assessment 9 \
\small (Answers are ordered exactly as in the question paper.)
\end{center}

\section*{Multiple Choice}
\begin{enumerate}[label=\arabic*.]
\item \textbf{Answer:} D
\\ \textit{Options:} A) 23.56 GHz; B) 50 MHz; C) 100 MHz; D) 42.58 MHz; E) 1 GHz
\\ \textit{Note:} The Larmor frequency for a proton in a magnetic field of 1 T is calculated using the formula f = γB, where γ (the gyromagnetic ratio for a proton) is approximately 42.58 MHz/T, resulting in a frequency of 42.58 MHz.
\item \textbf{Answer:} A
\\ \textit{Options:} A) 0.5 N; B) 0.2 N; C) 0.75 N; D) 0.1 N; E) 1 N
\\ \textit{Note:} The normality of the solution is 0.5 N because KOH is a strong base that dissociates to provide one hydroxide ion (OH⁻) per molecule, and using the formula for normality (N = equivalents of solute/volume of solution in liters), the calculation shows that 0.5 equivalents are present in 40 mL of the solution.
\item \textbf{Answer:} B
\\ \textit{Options:} A) 4 sigma and 1 pi; B) 3 sigma and 1 pi; C) 0 sigma and 4 pi; D) 3 sigma and 2 pi; E) 2 sigma and 1 pi
\\ \textit{Note:} Each resonance form of the nitrate ion, NO 3-, features three single bonds (sigma bonds) between the nitrogen and oxygen atoms, along with one double bond (which includes one sigma bond and one pi bond), resulting in a total of three sigma bonds and one pi bond across all resonance forms.
\item \textbf{Answer:} C
\\ \textit{Options:} A) 9.30 g O; B) 2.50 g O; C) 4.00 g O; D) 5.00 g O; E) 3.10 g O
\\ \textit{Note:} The correct answer is 4.00 g O because the stoichiometric ratio from the balanced reaction indicates that for every 2 atoms of phosphorus, 5 atoms of oxygen are needed, allowing us to calculate the required mass of oxygen based on the mass of phosphorus given.
\item \textbf{Answer:} A
\\ \textit{Options:} A) $4.20 \times 10^{-2}$ $\text{atm}$; B) $5.55 \times 10^{-2}$ $\text{atm}$; C) $6.33 \times 10^{-2}$ $\text{atm}$; D) $3.15 \times 10^{-2}$ $\text{atm}$; E) $8.21 \times 10^{-2}$ $\text{atm}$
\\ \textit{Note:} The correct answer is derived from applying the ideal gas law, \( PV = nRT \), where the number of moles \( n \) is calculated from the mass of neon, the volume \( V \) is given, and the temperature \( T \) is converted to Kelvin, leading to the calculated pressure of \( 4.20 \times 10^{-2} \text{atm} \).
\end{enumerate}

\section*{True / False}
\begin{enumerate}[label=\arabic*.]
\setcounter{enumi}{5}
\item \textbf{Answer:} True
\\ \textit{Options:} A) True; B) False
\\ \textit{Note:} The magnetogyric ratio of 23 Na is accurately determined to be 20.78 x $10^7$ T-1 s-1 based on its resonance frequency, which is a direct relationship in nuclear magnetic resonance (NMR) spectroscopy where the frequency is proportional to the magnetic field strength and the magnetogyric ratio.
\item \textbf{Answer:} True
\\ \textit{Options:} A) True; B) False
\\ \textit{Note:} When a solid melts, the transition from a highly ordered solid state to a more disordered liquid state results in an increase in entropy, and the energy required to overcome the intermolecular forces leads to a positive enthalpy change, making the statement true.
\item \textbf{Answer:} True
\\ \textit{Options:} A) True; B) False
\\ \textit{Note:} The statement is true because non-electrolytes typically have a molar weight around 100 g/mol, which corresponds to common organic compounds, allowing for convenient stoichiometric calculations in solution chemistry.
\item \textbf{Answer:} True
\\ \textit{Options:} A) True; B) False
\\ \textit{Note:} Hydrogen bonds involve a strong attraction between a hydrogen atom covalently bonded to a highly electronegative atom and another electronegative atom, making them stronger than dipole-dipole interactions, which arise from permanent dipoles, while induced dipole forces, or London dispersion forces, are the weakest due to their dependence on temporary dipoles that occur in all molecules.
\item \textbf{Answer:} True
\\ \textit{Options:} A) True; B) False
\\ \textit{Note:} The statement is true because the energy of one mole of photons can be calculated using the equation E = (hc/λ) multiplied by Avogadro's number, and for a wavelength of 300 nm, this calculation yields approximately 450 kJ/mol.
\end{enumerate}

\section*{Long Form Response}
\begin{enumerate}[label=\arabic*.]
\setcounter{enumi}{10}
\item \textbf{Answer:} Electrolysis in an electrolytic cell involves the decomposition of a compound, such as a zinc salt solution, through the application of an external electric current. During the electrolysis of zinc sulfate, zinc ions (\(Zn^{2+}\)) are reduced at the cathode, resulting in the deposition of metallic zinc. To calculate the equivalent weight of zinc, we first determine the total charge (Q) that passed through the cell using the formula \(Q = I \times t\), where \(I = 5.36 \, A\) and \(t = 30 \, min \times 60 \, s/min = 1800 \, s\). This results in \(Q = 5.36 \, A \times 1800 \, s = 9648 \, C\). Given that the deposition of 3.27 g of zinc corresponds to a transfer of 2 moles of electrons (since \(Zn^{2+} + 2 e^- \rightarrow Zn\)), the equivalent weight can be calculated using the formula \( \text{Equivalent weight} = \frac{\text{mass of zinc deposited}}{\text{number of equivalents}} \). Thus, the equivalent weight of zinc is determined to be \( \frac{3.27 \, g}{1.5} = 27.7 \, g\), confirming the equivalent weight based on the deposited mass.
\\ \textit{Note:} The answer correctly describes the electrolysis process by detailing the reduction of zinc ions at the cathode and accurately calculates the equivalent weight of zinc using the total charge derived from the current and time, demonstrating a solid understanding of Faraday's laws of electrolysis.
\item \textbf{Answer:} To analyze the electrochemical cell Pt/H₂/HCl (0.2 M)/HCl (0.8 M)/H₂/Pt at 25^\ensuremath{\circC}, we can apply the Nernst equation to calculate the cell voltage. Given that the standard electrode potential for the hydrogen electrode is 0 V, the cell voltage can be determined using the concentration difference of H⁺ ions in the two HCl solutions. The Nernst equation yields a cell voltage of 0.0006 volts. This low voltage is influenced by the ratio of the hydrogen ion concentrations in the two compartments, as well as temperature and pressure conditions, which dictate the electrochemical potential and the extent of the reaction.
\\ \textit{Note:} correct because it effectively applies the Nernst equation to account for the concentration difference of hydrogen ions between the two HCl solutions, leading to the calculated cell voltage of 0.0006 volts, which reflects the electrochemical behavior of the system at the specified conditions.
\end{enumerate}

\vfill
\noindent\textit{End of Answer Key}
\end{document}